%% kapa: exam=1 round=29 year=2018 period=1 form=A subject=civil_law subject_ko=민법 seq=16 number=16
%% source: Q-Net 감정평가사 기출문제 (cst00309) artlSeq=5202808 file=extracted/1008/5202808/2018년 제29회 감정평가사 1차시험 기출문제/2018년 제29회 감정평가사 1차시험 기출문제/1교시 A형.pdf
%% conversion: kordoc 4.24.51 -> tools/gen1.py (auto); review=proofread-sonnet; flags=none
\begin{question}{16}
소멸시효에 관한 설명으로 옳지 않은 것은? (다툼이 있으면 판례에 따름)
\begin{choices}
  \item 변제기가 도래한 단기소멸시효채권이 판결에 의해 확정된 경우 그 소멸시효는 5년으로 한다.
  \item 부작위를 목적으로 하는 채권의 소멸시효는 위반행위를 한 때로부터 진행한다.
  \item 최고는 6월내에 재판상의 청구, 파산절차참가, 화해를 위한 소환, 임의출석, 압류 또는 가압류, 가처분을 하지 아니하면 시효중단의 효력이 없다.
  \item 1년의 단기소멸시효에 걸리는 채권의 상대방이 그 채권의 발생원인이 된 계약에 기하여 가지는 반대채권은 특별한 사정이 없는 한 10년의 소멸시효에 걸린다.
  \item 소멸시효는 법률행위에 의하여 이를 배제, 연장 또는 가중할 수 없으나 이를 단축 또는 경감할 수 있다.
\end{choices}
\end{question}
\answer{1}
